\documentclass[11pt]{article}
\usepackage[margin=1in]{geometry}
\usepackage{amsmath,amssymb,amsthm,enumitem,hyperref}
\hypersetup{colorlinks=true,linkcolor=blue,urlcolor=blue,citecolor=blue}
\newtheorem{theorem}{Theorem}
\newcommand{\HH}{\mathcal H}
\title{An obstruction to the proposed binary Hankel product formula\\\large MathDB 375935}
\author{Principia Math In Progress}
\date{September 18, 2026}
\begin{document}
\maketitle
\begin{abstract}
The fixed-offset product formula in Sobolewski--Ulas's Conjecture 4.6
fails already at offset $2$. We evaluate that entire subsequence using
their determinant recurrence, then rule out every allowed shift and
prefactor. The obstruction persists even when polynomial prefactors
are replaced by arbitrary rational functions.
\end{abstract}

\section*{Task 1. Full disproof}
Let $s(a)$ be the number of ones in the binary expansion of $a$, and put
\[
 \HH(n)=\det\bigl[s(i+j)\bigr]_{0\le i,j<n},\qquad
 B_m=\prod_{j=1}^m(2^j-1)^{2^{m-j}},\qquad
 T_m=\sum_{j=1}^m\frac{2^{j-1}}{2^j-1}.
\]
Thus $\HH(n)=\HH(n,1)$ in the source's notation. The conjecture asserts
that, for each fixed $k\ge0$, there exist a fixed integer $r\ge0$, a sign
$\varepsilon$, a polynomial $P\in\mathbb Q[X]$, and a constant
$c\in\mathbb Q$ such that eventually
\[
 \HH(2^m+k)=\varepsilon P(2^m)(c+T_{m-r})B_{m-r}.                 \tag{1}
\]
The equality is required for all sufficiently large integers $m$.

\begin{theorem}
For every $m\ge2$, with $X=2^m$,
\[
 \boxed{\quad
 \HH(2^m+2)=\frac{B_m(2+T_{m-1})}{(2^{m-1}-1)(2^m-1)}
 =\frac{B_m}{A(X)}\bigl(T_m+d(X)\bigr),\quad}                  \tag{2}
\]
where
\[
 A(X)=(X/2-1)(X-1),\qquad d(X)=\frac32-\frac1{2(X-1)}.
\]
There are no $r\ge0$, $c\in\mathbb R$, and nonzero
$P\in\mathbb R(X)$ for which (1) holds eventually at $k=2$.
Consequently Conjecture 4.6 is false.
\end{theorem}

\subsection*{Exact evaluation}
We use the proved recurrence in \cite[Theorem 4.2 and Proposition 4.3]{SU}.
Here are its specializations, with all normalizations specified.
For $2^a<n\le2^{a+1}$ let $h(n,x)$ be the Hankel determinant of
the weighted digit sum with weights $(1,\ldots,1,x)$, where $x$
is the weight at position $a+1$. Let $g(n,x)$ be the determinant
of the $n\times n$ matrix whose first $n-1$ columns have entries
$s_{\mathbf w}(i+j+1)-s_{\mathbf w}(i+j)$ and whose last column is $1$.
For $m\ge2$ set
\[
 U_m(x)=h(2^m-1,x),\qquad V_m(x)=g(2^m-1,x).
\]
Direct $3\times3$ determinants give $U_2(x)=3-x$ and $V_2(x)=2-x$.
The source recurrence, applied on its second branch, gives for $m\ge3$
\[
 \begin{aligned}
 V_m(x)&=C_m(x)V_{m-1}(x/2),\\
 U_m(x)&=C_m(x)\left(U_{m-1}(x/2)+\frac{V_{m-1}(x/2)}{2-x}\right),\\
 C_m(x)&=2^{\,2^{m-1}-3}(2-x)^{2^{m-1}}.
 \end{aligned}                                                \tag{3}
\]
These identities follow with $n=2^m-1$ and $a=m-1$; in particular
$(-1)^n=-1$. Dividing the first two identities and iterating, initially
as identities of rational functions, gives
\[
 \frac{U_m(x)}{V_m(x)}
 =1+\sum_{j=0}^{m-2}\frac1{2-x/2^j}.                          \tag{4}
\]
Iteration of the first identity in (3) at $x=1$ also gives
\[
 \begin{aligned}
 V_m(1)
 &=\left(2-2^{-(m-2)}\right)
   \prod_{j=0}^{m-3}
    2^{\,2^{m-j-1}-3}\left(2-2^{-j}\right)^{2^{m-j-1}}\\
 &=(2^{m-1}-1)\prod_{j=1}^{m-2}(2^j-1)^{2^{m-j}}
   =\frac{B_m}{(2^{m-1}-1)(2^m-1)}.
 \end{aligned}                                                \tag{5}
\]
The product is empty at $m=2$. To check the powers of $2$ in the
second equality, their total exponent is
\[
 -(m-2)+\sum_{j=0}^{m-3}\bigl((1-j)2^{m-j-1}-3\bigr)=0.
\]
This is the elementary finite geometric-sum identity
$\sum_{j=0}^{m-3}(1-j)2^{m-j-1}=4m-8$.
Equation (4) at $x=1$ is $U_m(1)=V_m(1)(1+T_{m-1})$.

Finally use the source recurrence's first branch with $n=2^m+2$,
$a=m$, and all weights $1$. It is applicable for every $m\ge2$.
Its scalar prefactor and upper-right matrix entry are both $-1$,
so
\[
 \HH(2^m+2)=U_m(1)+V_m(1).
\]
Together with (4)--(5), this proves the first equality in (2).
The second follows from $T_m-T_{m-1}=X/(2(X-1))$.
For example, the values at $m=2,3,4,5$ are
$3,11,2403,1127461923$.

\subsection*{All fixed shifts and prefactors are impossible}
The positive constant
\[
 \lambda=\sum_{j=1}^{\infty}2^{-j}\log(2^j-1)
\]
is finite, since the summands are bounded by $j2^{-j}\log2$,
and it is positive already from $j=2$. Hence
\[
 \frac{\log B_m}{2^m}\longrightarrow\lambda.
\]
Also
\[
 T_m=\frac m2+\frac12\sum_{j=1}^m\frac1{2^j-1}
     =\frac m2+O(1).                                         \tag{6}
\]
By (2), $2^{-m}\log\HH(2^m+2)\to\lambda$.
If (1) holds with any nonzero rational function $P$, then
$\log|P(2^m)|=O(m)$ and $\log|c+T_{m-r}|=O(\log m)$.
Thus the same limit computed from (1) is $2^{-r}\lambda$.
It follows that $r=0$.

Absorb the sign into $P$ and define $R(X)=A(X)P(X)$.
Using (2), the remaining equality becomes
\[
 (R(2^m)-1)T_m=d(2^m)-cR(2^m).                              \tag{7}
\]
If $R\not\equiv1$, its difference from $1$ has only finitely many
zeros and poles. Consequently (7) would express $T_m$ eventually as
$F(2^m)$ for the fixed rational function
\[
 F(X)=\frac{d(X)-cR(X)}{R(X)-1}.
\]
This is impossible by (6): a rational function at infinity either
has a finite limit, or is asymptotic to $bX^e$ for some $b\ne0$
and positive integer $e$. Neither behavior at $X=2^m$ is
asymptotic to $m/2$. If $R\equiv1$, (7) instead forces
$c=d(2^m)$ for every sufficiently large $m$, also impossible because
$d(2^m)$ is strictly increasing. This proves the theorem.

The argument concerns the exact eventual identity in the source.
It does not assert that all asymptotic approximations by related
products fail. In fact, (2) supplies the correct exact formula at
offset $2$, with a variable rational additive term.

\section*{Verification and scope}
The accompanying script checks direct rational Hankel determinants,
both recurrence branches against their defining matrices, the
fixed-offset formula at larger indices, and the finite identities
used in the derivation. These finite checks support the computation;
the proof excluding every $r,P,c$ is the argument above.
The determinant recurrences and the special families at offsets
$-1,0,1$ belong to Sobolewski and Ulas. The new result here is the
offset-$2$ evaluation and its obstruction to Conjecture 4.6.
Review is by the authoring agent only, without independent peer review.

\section*{Tasks 2--5. Follow-up directions}
\begin{enumerate}[start=2,leftmargin=*]
\item Determine whether every fixed offset admits an eventual formula
$\HH(2^m+k)=B_m(A_k(2^m)T_m+D_k(2^m))$ with rational functions
$A_k,D_k$, and construct them from the binary digits of $k$.
\item Classify offsets for which $D_k/A_k$ is constant, recovering a
single factor $T_m+c_k$ with a rational prefactor.
\item Derive complete asymptotic expansions of the offset families,
including the powers of $2^{-m}$ and the terms proportional to $m$.
\item Extend the fixed-offset analysis to the polynomial weights $w_j=t^j$
and identify how the exceptional values $t=2\zeta$ change it.
\end{enumerate}

\begin{thebibliography}{9}
\raggedright
\bibitem{SU} B. Sobolewski and M. Ulas,
\emph{Hankel determinants of weighted binary sums of digits},
\href{https://arxiv.org/abs/2607.09376v1}{arXiv:2607.09376v1}, July 10, 2026.
Theorem 4.2, Proposition 4.3 and Conjecture 4.6.
\bibitem{DB} MathDB,
\href{https://mathdb.com/p/375935}{Problem 375935}, accessed September 18, 2026.
\end{thebibliography}
\end{document}
